\documentclass[journal]{IEEEtran}

\usepackage{amsmath,amsfonts}
\usepackage{amssymb}
\usepackage{graphicx}
\usepackage{cite}
\usepackage{booktabs}
\usepackage{empheq}

\begin{document}

\title{Distributed Model Predictive Control for Energy Management of a PV--Battery--Hydrogen DC Microgrid}

\author{Amit~Uniyal}

\maketitle

%-------------------------------------------------
\begin{abstract}
This paper presents a distributed model predictive control (DMPC) framework
for energy management in a DC microgrid integrating solar photovoltaic (PV)
generation, battery energy storage, and hydrogen-based energy systems. The
proposed system incorporates a 36\,kW PV array, a battery storage unit, a
10\,kW proton-exchange membrane fuel cell (PEMFC), and a 10\,kW alkaline
electrolyzer (AE) interconnected through a common 300\,V DC bus. A
distributed control strategy is developed to ensure coordinated operation of
subsystems under varying load and generation conditions. The control
objective includes DC bus voltage regulation, power balance, and optimal
energy utilization. Simulation results demonstrate improved stability,
efficient energy distribution, and effective handling of renewable
intermittency.
\end{abstract}

\begin{IEEEkeywords}
DC microgrid, Distributed MPC, Solar PV, Hydrogen energy, Fuel cell,
Electrolyzer, Energy management, Grey Wolf Optimizer.
\end{IEEEkeywords}

%-------------------------------------------------
\section{Introduction}

The increasing penetration of solar photovoltaic (PV) systems in modern
power systems necessitates advanced energy management strategies to handle
intermittency and ensure reliable operation. Hybrid energy storage systems
combining batteries and hydrogen technologies provide a promising solution
for both short-term and long-term energy balancing.

Conventional rule-based energy management approaches suffer from limited
adaptability under dynamic conditions. Model predictive control (MPC),
particularly in distributed form, offers improved flexibility and
scalability by allowing each subsystem to optimize its own local objective
while exchanging only a minimal set of communication variables with
neighbouring controllers.

This paper proposes a distributed economic MPC (DEMPC)-based energy
management system for a DC microgrid integrating PV, battery, PEMFC, and
alkaline electrolyzer. The main contributions are:
\begin{itemize}
\item Development of a detailed nonlinear dynamic model of the hybrid
      PV--battery--hydrogen DC microgrid at the converter level.
\item Formulation of a distributed economic MPC strategy with subsystem-specific
      droop-based cost functions and a supervisory energy management system (EMS).
\item Validation under dynamic load and irradiance scenarios demonstrating
      DC bus voltage regulation and coordinated hydrogen system operation.
\end{itemize}

%-------------------------------------------------
\section{System Description}

The DC microgrid consists of a PV array, a battery energy storage system
(BESS), a PEMFC stack, an alkaline electrolyzer, and a variable DC load
connected through a common DC bus of nominal voltage $V_{dc}^{*} = 300$\,V.
Each source and storage element is interfaced via a dedicated DC--DC
converter. The overall system power balance is:

\begin{equation}
P_{pv} + P_{fc} + P_{bat} = P_{load} + P_{ae} + P_{loss}
\label{eq:power_balance}
\end{equation}

where $P_{pv}$, $P_{fc}$, $P_{bat}$, $P_{ae}$, and $P_{load}$ denote the
PV, fuel cell, battery, electrolyzer, and load powers, respectively, and
$P_{loss}$ accounts for converter losses (neglected in the ideal model).

%-------------------------------------------------
\section{Component Modeling}
\label{sec:modeling}

This section presents the dynamic modeling of all major components. The
system is modeled at the converter level to capture switching dynamics and
nonlinear behavior, which is essential for the predictive control
implementation.

%-------------------------------------------------
\subsection{Photovoltaic Array}
\label{subsec:pv}

The PV array consists of $N_s = 5$ modules in series and $N_{sh} = 20$
parallel strings. Each module is modeled using the single-diode equivalent
circuit with both series resistance $R_s$ and shunt resistance $R_{sh}$.
For an array of $N_s \times N_{sh}$ modules the implicit current equation is:

\begin{equation}
\begin{split}
I_{pv} = \; & N_{sh}\,I_{ph} - N_{sh}\,I_0\!\left(
               \exp\!\left(\frac{V_{pv}+I_{pv}R_{s,eq}}{aN_sV_t}\right)-1
             \right) \\
           &- \frac{V_{pv}+I_{pv}R_{s,eq}}{R_{sh,eq}}
\end{split}
\label{eq:pv_iv}
\end{equation}

where the array-equivalent series and shunt resistances are:
\begin{equation}
R_{s,eq} = \frac{N_s}{N_{sh}}R_s, \qquad
R_{sh,eq} = \frac{N_s}{N_{sh}}R_{sh}
\label{eq:pv_req}
\end{equation}

and $V_t = kT/q$ is the cell thermal voltage at operating temperature $T$,
$a$ is the diode ideality factor, and $N_{cell}$ is the number of cells per
module. Because \eqref{eq:pv_iv} is implicit in $I_{pv}$, it is solved at
each time step using Newton--Raphson iteration.

The photocurrent depends on irradiance $G$ and temperature $T$:
\begin{equation}
I_{ph} = \frac{G}{1000}\!\left(I_{scn} + k_i\,(T - T_n)\right)
\label{eq:pv_iph}
\end{equation}

where $I_{scn}$ is the short-circuit current at standard test conditions
(STC: $G_n = 1000$\,W/m$^2$, $T_n = 25\,^\circ$C) and $k_i$ is the
temperature coefficient.

The saturation current at operating temperature is:
\begin{equation}
I_0 = I_{on}\!\left(\frac{T}{T_n}\right)^{\!3}
      \exp\!\left[\frac{qE_{gap}}{ka}\!\left(\frac{1}{T_n}-\frac{1}{T}\right)\right]
\label{eq:pv_io}
\end{equation}

where the nominal saturation current at STC is:
\begin{equation}
I_{on} = \frac{I_{scn}}{\exp\!\left(\dfrac{V_{ocn}}{aV_{t,n}}\right)-1}
\label{eq:pv_ion}
\end{equation}

The PV array output power is:
\begin{equation}
P_{pv} = V_{pv}\,I_{pv}
\label{eq:pv_power}
\end{equation}

%-------------------------------------------------
\subsection{PV Boost Converter}
\label{subsec:pv_conv}

The PV array is interfaced with the DC bus through a unidirectional boost
converter. The converter introduces two dynamic states: the PV-side
capacitor voltage $v_{pv}$ and the boost inductor current $i_L$. Their
state equations are:

\begin{equation}
C_s\,\frac{dv_{pv}}{dt} = i_{pv}(v_{pv}) - i_L
\label{eq:pv_cap}
\end{equation}

\begin{equation}
L_s\,\frac{di_L}{dt} = v_{pv} - (1-d_s)\,V_{dc}
\label{eq:pv_ind}
\end{equation}

The current injected into the DC bus is:
\begin{equation}
i_s = (1-d_s)\,i_L
\label{eq:pv_is}
\end{equation}

where $d_s \in [0,1]$ is the boost converter duty cycle and $C_s$, $L_s$
are the filter capacitance and inductance, respectively.

%-------------------------------------------------
\subsection{Alkaline Electrolyzer Stack}
\label{subsec:ae_stack}

The alkaline electrolyzer (AE) stack consists of $N_{ae} = 30$ cells in
series. The per-cell terminal voltage combines the reversible voltage, an
ohmic drop, and a Tafel overvoltage:

\begin{equation}
\begin{split}
V_{ae}(i_{ae}) = \; & N_{ae}\!\left[\vphantom{\frac{t_1}{A}}\right. U_{rev} + \frac{r_1 +   r_2\,T_{ae}}{A}\,i_{ae} \\
                    &+ t_e\ln\!\left(
                        \frac{t_1+t_2/T_{ae}+t_3/T_{ae}^2}{A}\,i_{ae}+1
                      \right)
                    \left.\vphantom{\frac{t_1}{A}}\right]
\end{split}
\label{eq:ae_voltage}
\end{equation}

where $U_{rev} = 1.229$\,V is the reversible cell voltage, $r_1$, $r_2$
are ohmic resistance coefficients, $t_e$, $t_1$, $t_2$, $t_3$ are Tafel
overvoltage parameters, $A$ is the active cell area, and $T_{ae}$ is the
cell temperature.

The AE stack power consumed is:
\begin{equation}
P_{ae} = V_{ae}(i_{ae})\,i_{ae}
\label{eq:ae_power}
\end{equation}

The Faraday (current) efficiency, which accounts for parasitic current
losses, is modeled as:
\begin{equation}
\eta_F = \frac{(i_{ae}/A)^2}{f_1 + (i_{ae}/A)^2}\,f_2
\label{eq:ae_etaF}
\end{equation}

where $f_1$ and $f_2$ are empirical efficiency parameters. The hydrogen
production rate is then:
\begin{equation}
\dot{N}_{H_2} = \frac{\eta_F\,N_{ae}\,i_{ae}}{2F}
\label{eq:ae_nh2}
\end{equation}

where $F = 96485.3$\,C/mol is the Faraday constant.

%-------------------------------------------------
\subsection{AE Buck Converter}
\label{subsec:ae_conv}

The electrolyzer is supplied from the DC bus through a unidirectional buck
converter. The inductor current dynamics are:

\begin{equation}
L_{ae}\,\frac{di_{ae}}{dt} = d_{ae}\,V_{dc} - V_{ae}(i_{ae})
\label{eq:ae_conv}
\end{equation}

The current drawn from the DC bus is:
\begin{equation}
i_a = d_{ae}\,i_{ae}
\label{eq:ae_ia}
\end{equation}

where $d_{ae} \in [0,1]$ is the buck converter duty cycle and $L_{ae}$ is
the converter inductance.

%-------------------------------------------------
\subsection{PEMFC Stack}
\label{subsec:fc_stack}

The proton-exchange membrane fuel cell (PEMFC) stack consists of $N_{pe} = 300$
cells in series operating at temperature $T_{pe} = 323.15$\,K. The stack
output voltage is composed of the thermodynamic Nernst voltage less three
polarization losses:

\begin{equation}
V_{pe}(i_{pe}) = N_{pe}\!\left(E_{nernst} + v_{act} + v_{ohm} + v_{con}\right)
\label{eq:fc_voltage}
\end{equation}

\subsubsection*{Nernst EMF}
The open-circuit thermodynamic voltage per cell is:
\begin{equation}
\begin{split}
    E_{nernst} = \; & 1.229 - 8.45\!\times\!10^{-4}(T_{pe}-298.15)\\
             & + 4.3085\!\times\!10^{-5}\,T_{pe}\!\left(\ln P_{H_2}+0.5\ln P_{O_2}\right)    
\end{split}
\label{eq:fc_nernst}
\end{equation}

where $P_{H_2}$ and $P_{O_2}$ are the hydrogen and oxygen partial pressures
in atm.

\subsubsection*{Activation polarization}
\begin{equation}
\begin{split}
v_{act} = \; & -0.9514 + 3.12\!\times\!10^{-3}T_{pe}
          + 7.4\!\times\!10^{-5}T_{pe}\ln C_{O_2} \\
          &- 1.87\!\times\!10^{-4}T_{pe}\ln i_{pe}    
\end{split}
\label{eq:fc_vact}
\end{equation}

where $C_{O_2}$ is the oxygen concentration at the cathode--catalyst
interface:
\begin{equation}
C_{O_2} = \frac{P_{O_2}}{5.08\!\times\!10^6\,\exp(-498/T_{pe})}
\label{eq:fc_co2}
\end{equation}

\subsubsection*{Ohmic polarization}
The Nafion membrane resistivity is:
\begin{equation}
r_m = \frac{181.6\!\left[1 + 0.03J + 0.062(T_{pe}/303)^2 J^{2.5}\right]}
           {(\lambda - 0.634 - 3J)\exp\!\left[4.18(T_{pe}-303)/T_{pe}\right]}
\label{eq:fc_rm}
\end{equation}

where $J = i_{pe}/A_{pe}$ is the current density (A/cm$^2$) and $\lambda$
is the membrane water content. The ohmic loss per cell is:
\begin{equation}
v_{ohm} = -i_{pe}\!\left(\frac{r_m\,l_m}{A_{pe}} + R_c\right)
\label{eq:fc_vohm}
\end{equation}

where $l_m$ is the membrane thickness and $R_c$ is the contact resistance.

\subsubsection*{Concentration polarization}
\begin{equation}
v_{con} = 0.016\ln\!\left(1 - \frac{i_{pe}}{A_{pe}\,J_{max}}\right)
\label{eq:fc_vcon}
\end{equation}

where $J_{max}$ is the maximum current density.

\subsubsection*{Stack power and hydrogen consumption}
\begin{equation}
P_{pe} = V_{pe}(i_{pe})\,i_{pe}
\label{eq:fc_power}
\end{equation}

\begin{equation}
\dot{q}_{H_2} = \frac{N_{pe}\,i_{pe}}{2F}
\label{eq:fc_qh2}
\end{equation}

%-------------------------------------------------
\subsection{PEMFC Boost Converter}
\label{subsec:fc_conv}

The fuel cell is connected to the DC bus through a unidirectional boost
converter. The inductor current dynamics are:

\begin{equation}
L_{pe}\,\frac{di_{pe}}{dt} = V_{pe}(i_{pe}) - (1-d_{pe})\,V_{dc}
\label{eq:fc_conv}
\end{equation}

The current injected into the DC bus is:
\begin{equation}
i_p = (1-d_{pe})\,i_{pe}
\label{eq:fc_ip}
\end{equation}

where $d_{pe} \in [0,1]$ is the boost converter duty cycle.

%-------------------------------------------------
\subsection{Battery Energy Storage System}
\label{subsec:bess}

The battery is modeled using a linear internal resistance model. The
terminal voltage is:
\begin{equation}
V_{bat} = V_{nom} - R_{int}\,I_{bat}
\label{eq:bat_v}
\end{equation}

where $V_{nom} = 240$\,V is the nominal voltage and $R_{int}$ is the
internal resistance. The battery output power is:
\begin{equation}
P_{bat} = V_{bat}\,I_{bat}
\label{eq:bat_p}
\end{equation}

The state-of-charge dynamics follow from Coulomb counting:
\begin{equation}
\frac{d\,\text{SOC}}{dt} = -\frac{I_{bat}}{C_{bat}}
\label{eq:bat_soc}
\end{equation}

where $C_{bat} = C_{Ah} \times 3600$\,As is the battery capacity
($C_{Ah} = 100$\,Ah). The SOC is constrained to safe limits:
\begin{equation}
\text{SOC}_{min} \leq \text{SOC} \leq \text{SOC}_{max}
\label{eq:bat_soc_lim}
\end{equation}

%-------------------------------------------------
\subsection{Battery Bidirectional Converter}
\label{subsec:bess_conv}

The battery is interfaced through a bidirectional DC--DC buck-boost
converter. The inductor current dynamics are:
\begin{equation}
L_b\,\frac{di_b}{dt} = V_{bat} - (1-d_b)\,V_{dc}
\label{eq:bess_conv}
\end{equation}

The current injected into (or drawn from) the DC bus is:
\begin{equation}
i_{b,dc} = (1-d_b)\,i_b
\label{eq:bess_ibus}
\end{equation}

where $d_b \in [0,1]$ is the converter duty cycle. Positive $i_b$ denotes
discharging; negative $i_b$ denotes charging.

%-------------------------------------------------
\subsection{DC Bus Dynamics}
\label{subsec:dcbus}

The DC bus voltage dynamics are governed by the capacitor KCL:
\begin{equation}
C_{dc}\,\frac{dV_{dc}}{dt} = i_s + i_p + i_{b,dc} - i_a - i_l
\label{eq:dcbus}
\end{equation}

The load is modeled as a constant power load (CPL), giving a nonlinear
current term:
\begin{equation}
i_l = \frac{P_{load}}{V_{dc}}
\label{eq:dcbus_load}
\end{equation}

%-------------------------------------------------
\subsection{Complete State-Space Representation}
\label{subsec:statespace}

Combining all subsystem models, the full system state vector is:
\begin{equation}
x = \begin{bmatrix}
v_{pv} & i_L & i_{ae} & i_{pe} & V_{dc} & i_b & \text{SOC}
\end{bmatrix}^T \in \mathbb{R}^7
\label{eq:state_vec}
\end{equation}

The control input vector (converter duty cycles) is:
\begin{equation}
u = \begin{bmatrix}d_s & d_{ae} & d_{pe} & d_b\end{bmatrix}^T \in [0,1]^4
\label{eq:ctrl_vec}
\end{equation}

and the disturbance vector is:
\begin{equation}
w = \begin{bmatrix}G & T & P_{load}\end{bmatrix}^T
\label{eq:dist_vec}
\end{equation}

The complete nonlinear state equations are:
\begin{equation}
\begin{aligned}
C_s\,\dot{v}_{pv}      &= i_{pv}(v_{pv},G,T) - i_L \\
L_s\,\dot{i}_L         &= v_{pv} - (1-d_s)\,V_{dc} \\
L_{ae}\,\dot{i}_{ae}   &= d_{ae}\,V_{dc} - V_{ae}(i_{ae}) \\
L_{pe}\,\dot{i}_{pe}   &= V_{pe}(i_{pe}) - (1-d_{pe})\,V_{dc} \\
C_{dc}\,\dot{V}_{dc}   &= i_s + i_p + i_{b,dc} - i_a - P_{load}/V_{dc} \\
L_b\,\dot{i}_b         &= V_{bat} - (1-d_b)\,V_{dc} \\
\dot{\text{SOC}}       &= -I_{bat}/C_{bat}
\end{aligned}
\label{eq:full_statespace}
\end{equation}

This nonlinear system forms the prediction model used in the DEMPC
formulation described in Section~\ref{sec:dempc}.

%-------------------------------------------------
%-------------------------------------------------
\section{Energy Management System}
\label{sec:ems}

A supervisory energy management system (EMS) operates at the highest level
of the control hierarchy and determines which hydrogen subsystems are
permitted to operate at each sampling instant. The EMS acts as a
discrete-event coordinator: it does not compute duty cycles directly, but
rather issues binary activation flags that enable or disable the alkaline
electrolyzer (AE) and the PEMFC via virtual contactors. This architecture
decouples the slow energy-scheduling decisions from the fast converter-level
MPC loops.

%-------------------------------------------------
\subsection{Net Power Balance}
\label{subsec:ems_pnet}

At every sampling instant $k$, the EMS evaluates the available power
balance using the PV maximum power $P_{pv,max}(k)$, computed by sweeping
the PV array model across its voltage range and extracting the global peak:

\begin{equation}
P_{net}(k) = P_{pv,max}(k) - P_{load}(k)
\label{eq:ems_pnet}
\end{equation}

A positive value of $P_{net}$ indicates a generation surplus that must be
absorbed; a negative value indicates a generation deficit that must be
compensated. The battery handles imbalances within its power rating
$P_{bat}^{max}$; beyond this limit, or when the battery SOC reaches a
boundary, the hydrogen subsystems are engaged.

%-------------------------------------------------
\subsection{Activation Logic}
\label{subsec:ems_logic}

The EMS generates two binary activation flags,
$\zeta_a \in \{0,1\}$ for the electrolyzer and
$\zeta_p \in \{0,1\}$ for the fuel cell, according to the following
priority-ordered logic:

\textbf{Step 1 — Default state.} Both flags are initialised to zero at
every time step:
\begin{equation}
\zeta_a(k) = 0, \qquad \zeta_p(k) = 0
\label{eq:ems_default}
\end{equation}

\textbf{Step 2 — Electrolyzer activation.} The electrolyzer is enabled
($\zeta_a = 1$) if either of the following surplus conditions holds:

\begin{equation}
\begin{split}
\zeta_a = 1 \quad \text{if} \quad
\begin{aligned}
& P_{net} > P_{bat}^{max} \\
& \text{or} \\
& \bigl(P_{net} > 0 \;\wedge\; \mathrm{SOC} \geq \mathrm{SOC}_{max}\bigr)
\end{aligned}
\label{eq:ems_ae}
\end{split}
\end{equation}

The first condition handles cases in which the generation surplus exceeds
the battery charging power limit and would otherwise cause the DC bus
voltage to rise uncontrollably. The second condition prevents the battery
from being overcharged when it is already near full capacity but a surplus
still exists.

\textbf{Step 3 — Fuel cell activation.} The fuel cell is enabled
($\zeta_p = 1$) if either of the following deficit conditions holds:

\begin{equation}
\zeta_p = 1 \quad \text{if} \quad
\begin{aligned}
& P_{net} < -P_{bat}^{max} \\
& \text{or} \\
& \bigl(P_{net} < 0 \;\wedge\; \mathrm{SOC} \leq \mathrm{SOC}_{min}\bigr)
\end{aligned}
\label{eq:ems_fc}
\end{equation}

The first condition handles cases in which the deficit exceeds the battery
discharge power limit. The second condition prevents deep discharge when the
battery SOC has reached its lower bound.

The numerical values of all EMS thresholds are summarised in
Table~\ref{tab:ems_params}.

\begin{table}[!h]
\caption{EMS Threshold Parameters}
\label{tab:ems_params}
\centering
\begin{tabular}{@{}lll@{}}
\toprule
Parameter & Symbol & Value \\
\midrule
Battery power limit    & $P_{bat}^{max}$    & 5\,kW \\
Minimum SOC            & $\mathrm{SOC}_{min}$ & 0.20  \\
Maximum SOC            & $\mathrm{SOC}_{max}$ & 0.90  \\
\bottomrule
\end{tabular}
\end{table}

%-------------------------------------------------
\subsection{Mutual Exclusivity and Priority}
\label{subsec:ems_mutex}

The two activation conditions \eqref{eq:ems_ae} and \eqref{eq:ems_fc} are
evaluated independently, which means that in pathological operating
conditions — for example, a large surplus coinciding with a depleted battery
— both $\zeta_a$ and $\zeta_p$ could simultaneously equal one. In normal
operation this is physically inconsistent: it would require the electrolyzer
to absorb power while the fuel cell simultaneously produces power from the
same hydrogen store.

To prevent simultaneous activation, an explicit priority rule is enforced
after Steps~2 and~3:

\begin{equation}
\text{if } \zeta_a = 1 \;\Rightarrow\; \zeta_p = 0
\label{eq:ems_priority}
\end{equation}

Electrolyzer absorption is given priority over fuel-cell dispatch. This
choice reflects the physical preference to store surplus energy as hydrogen
rather than to simultaneously consume and produce it. Under this rule, the
feasible flag combinations and their physical interpretations are:

\begin{equation}
(\zeta_a,\, \zeta_p) \in
\bigl\{(0,0),\;(1,0),\;(0,1)\bigr\}
\label{eq:ems_states}
\end{equation}

corresponding to \emph{idle}, \emph{electrolyzer active}, and
\emph{fuel cell active} modes, respectively.

%-------------------------------------------------
\subsection{Virtual Contactor Enforcement}
\label{subsec:ems_contactor}

When a subsystem is deactivated ($\zeta_a = 0$ or $\zeta_p = 0$), a
virtual contactor is opened: the corresponding converter duty cycle is set
to zero and the associated inductor current state is forced to zero at the
end of every integration step:

\begin{equation}
\zeta_a = 0 \;\Rightarrow\; d_{ae} = 0,\quad i_{ae} \leftarrow 0
\label{eq:ems_contactor_ae}
\end{equation}
\begin{equation}
\zeta_p = 0 \;\Rightarrow\; d_{pe} = 0,\quad i_{pe} \leftarrow 0
\label{eq:ems_contactor_fc}
\end{equation}

This hard clamping avoids spurious current dynamics in the inductor states
of inactive converters and ensures that the DEMPC prediction models of the
active subsystems receive consistent communicated values (zero power and
zero bus current) from inactive subsystems.

The complete EMS decision procedure executed at each sampling instant is
summarised in Algorithm~\ref{alg:ems}.

\begin{figure}[!h]
\centering
\begin{minipage}{0.9\linewidth}
\renewcommand{\arraystretch}{1.2}
\hrule\vspace{4pt}
\textbf{Algorithm~1:} Supervisory EMS Decision Procedure
\vspace{4pt}\hrule\vspace{4pt}
\begin{enumerate}
  \item Compute $P_{pv,max}(k)$ via MPP sweep of PV array model.
  \item Evaluate $P_{net}(k) = P_{pv,max}(k) - P_{load}(k)$.
  \item Set $\zeta_a \leftarrow 0$, $\zeta_p \leftarrow 0$.
  \item \textbf{If} $P_{net} > P_{bat}^{max}$ \textbf{or}
        $(P_{net} > 0$ \textbf{and} $\mathrm{SOC} \geq \mathrm{SOC}_{max})$:
        set $\zeta_a \leftarrow 1$.
  \item \textbf{If} $P_{net} < -P_{bat}^{max}$ \textbf{or}
        $(P_{net} < 0$ \textbf{and} $\mathrm{SOC} \leq \mathrm{SOC}_{min})$:
        set $\zeta_p \leftarrow 1$.
  \item \textbf{If} $\zeta_a = 1$: set $\zeta_p \leftarrow 0$
        (electrolyzer priority rule).
  \item Pass $(\zeta_a, \zeta_p)$ to local DEMPC controllers and
        enforce virtual contactors~\eqref{eq:ems_contactor_ae}--\eqref{eq:ems_contactor_fc}.
\end{enumerate}
\vspace{4pt}\hrule
\end{minipage}
\label{alg:ems}
\end{figure}

%-----------------------------------------------------------------------------------
\section{Distributed Economic Model Predictive Control}
\label{sec:dempc}

\subsection{Architecture Overview}

The DEMPC framework decomposes the centralized problem into four
independent local controllers: PV subsystem (PVS), AE subsystem (AES),
PEMFC subsystem (PEMFCS), and BESS. At each sampling instant $T_s$, each
local controller:
\begin{enumerate}
\item Receives communicated power and current signals from the other
      subsystems (from the previous time step).
\item Solves its own finite-horizon optimal control problem over a
      prediction horizon $N_p$.
\item Applies the first element of the optimal duty-cycle sequence.
\item Broadcasts its power and bus current to the other controllers.
\end{enumerate}

The DC bus voltage reference is updated with a reference smoothing filter:
\begin{equation}
V_{dc,ref}(k+1) = V_{dc,ref}(k) + \frac{1}{N_s}\!\left(V_{dc}^* - V_{dc,ref}(k)\right)
\label{eq:vdc_ref}
\end{equation}

where $N_s = 10$ is the smoothing horizon and $V_{dc}^* = 300$\,V.

%-------------------------------------------------
\subsection{Local Cost Function — PV Subsystem (PVS)}

The PVS controller minimizes a cost that tracks the maximum power point
(MPPT), regulates the DC bus voltage, and tracks the MPPT inductor current:

\begin{equation}
J_{pvs} = \sum_{l=1}^{N_p}\!\Bigl[
  w_{s2}\,\ell_{s2}^l + w_{s3}\,\ell_{s3}^l + w_{s4}\,\ell_{s4}^l
\Bigr]
\label{eq:cost_pvs}
\end{equation}

where the normalized stage costs at prediction step $l$ are:
\begin{empheq}[left=\empheqlbrace]{align}
\ell_{s2}^l &= \left(\frac{P_{pv}^l - P_{max}}{\max(P_{max},1)}\right)^{\!2}
\label{eq:pvs_l2}\\
\ell_{s3}^l &= \left(\frac{V_{dc}^l - V_{dc,ref}}{V_{dc,ref}}\right)^{\!2}
\label{eq:pvs_l3}\\
\ell_{s4}^l &= \left(\frac{i_L^l - i_{L,ref}}{\max(i_{L,ref},1)}\right)^{\!2}
\label{eq:pvs_l4}
\end{empheq}

with weights $w_{s2} = 800$, $w_{s3} = 20000$, $w_{s4} = 400$.
The inductor current reference $i_{L,ref}$ is set to the MPP current
$I_{mpp}$ computed from a 500-point voltage sweep.

%-------------------------------------------------
\subsection{Local Cost Function — AE Subsystem (AES)}

The AES controller drives the electrolyzer power toward a surplus
absorption target while regulating the bus voltage:

\begin{equation}
J_{aes} = \sum_{l=1}^{N_p}\!\Bigl[
  w_{a1}\,\ell_{a1}^l + w_{a2}\,\ell_{a2}^l
\Bigr]
\label{eq:cost_aes}
\end{equation}

\begin{empheq}[left=\empheqlbrace]{align}
    \ell_{a1}^l &= \left(\frac{P_{ae}^l - P_{tgt,a}}{\max(P_{tgt,a},1)}\right)^{\!2}
    \label{eq:aes_l1}\\
    \ell_{a2}^l &= \left(\frac{V_{dc}^l - V_{dc,ref}}{V_{dc,ref}}\right)^{\!2}
    \label{eq:aes_l2}
\end{empheq}

with weights $w_{a1} = 400$, $w_{a2} = 20000$. The surplus absorption
target incorporates a voltage-droop correction:
\begin{equation}
P_{tgt,a} = \max\!\Bigl(
  P_{pv} + P_{fc} - P_{load} + K_{vdc}(V_{dc} - V_{dc,ref}),\; 0
\Bigr)
\label{eq:aes_target}
\end{equation}

where $K_{vdc} = 400$\,W/V is the droop gain.

%-------------------------------------------------
\subsection{Local Cost Function — PEMFC Subsystem (PEMFCS)}

The PEMFCS controller drives the fuel cell toward a deficit compensation
target:

\begin{equation}
J_{pemfcs} = \sum_{l=1}^{N_p}\!\Bigl[
  w_{p1}\,\ell_{p1}^l + w_{p2}\,\ell_{p2}^l
\Bigr]
\label{eq:cost_fc}
\end{equation}

\begin{empheq}[left=\empheqlbrace]{align}
\ell_{p1}^l &= \left(\frac{P_{pe}^l - P_{tgt,p}}{\max(P_{tgt,p},1)}\right)^{\!2}
\label{eq:fc_l1}\\
\ell_{p2}^l &= \left(\frac{V_{dc}^l - V_{dc,ref}}{V_{dc,ref}}\right)^{\!2}
\label{eq:fc_l2}
\end{empheq}

with weights $w_{p1} = 400$, $w_{p2} = 20000$. The deficit compensation
target is:
\begin{equation}
P_{tgt,p} = \max\!\Bigl(
  P_{load} + P_{ae} - P_{pv} - K_{vdc}(V_{dc} - V_{dc,ref}),\; 0
\Bigr)
\label{eq:fc_target}
\end{equation}

%-------------------------------------------------
\subsection{Local Cost Function — BESS}

The BESS controller minimizes a composite cost penalizing voltage
deviation, power imbalance, duty-cycle variation, and constraint
violations:

\begin{equation}
J_{bess} = \sum_{l=1}^{N_p}\!\Bigl[
  w_v\,\ell_v^l + w_{du}\,\ell_{du}^l + w_p\,\ell_p^l
  + w_{lim}\,\ell_{lim}^l + w_{soc}\,\ell_{soc}^l
\Bigr]
\label{eq:cost_bess}
\end{equation}

\begin{empheq}[left=\empheqlbrace]{align}
\ell_v^l   &= \left(\frac{V_{dc}^l - V_{dc,ref}}{V_{dc,ref}}\right)^{\!2} \\
\ell_{du}^l &= (d_b^l - d_b^{l-1})^2 \\
\ell_p^l   &= \left(\frac{P_{pv}+P_{pe}+P_{bat}^l-P_{ae}-P_{load}}
                         {\max(P_{load},1)}\right)^{\!2} \\
\ell_{lim}^l &= \left(\frac{\max(0,\,|P_{bat}^l|-P_{bat}^{max})}
                            {P_{bat}^{max}}\right)^{\!2} \\
\ell_{soc}^l &= \left(\frac{\Psi_{soc}(P_{bat}^l,\,\text{SOC}^l)}
                             {P_{bat}^{max}}\right)^{\!2}
\label{eq:bess_costs}
\end{empheq}

with weights $w_v = 20000$, $w_{du} = 100$, $w_p = 5000$,
$w_{lim} = w_{soc} = 50000$, and $P_{bat}^{max} = 5$\,kW. The SOC
violation penalty is:
\begin{equation}
\Psi_{soc} = \begin{cases}
|P_{bat}^l| & \text{if } \text{SOC}^l \geq 0.9 \text{ and } P_{bat}^l < 0 \\
|P_{bat}^l| & \text{if } \text{SOC}^l \leq 0.2 \text{ and } P_{bat}^l > 0 \\
0           & \text{otherwise}
\end{cases}
\label{eq:soc_penalty}
\end{equation}

%-------------------------------------------------
\subsection{Optimization and Prediction}

At each sampling instant, each local controller optimizes over a sequence
of $N_p$ duty cycles. The local prediction model integrates the subsystem
state equations using a forward-Euler scheme with four sub-steps per
sampling interval $T_s = 40\,\mu$s. The duty cycles of all other
subsystems are frozen at their communicated values during the local
prediction (distributed approximation).

Each optimization problem takes the form:
\begin{equation}
\min_{d^{[1:N_p]}} \; J_i(d^{[1:N_p]})
\quad \text{s.t.} \quad d^l \in [\underline{d}_i,\, \overline{d}_i], \;\forall l
\label{eq:local_opt}
\end{equation}

For the baseline controller, the feasible duty-cycle space is discretized
into a 7-point adaptive grid centered on the previous duty cycle, and all
$7^{N_p}$ sequences are enumerated. For the Grey Wolf Optimizer (GWO)
variant, the same cost function is minimized using a population-based
metaheuristic, reducing computational complexity while retaining solution
quality.

%-------------------------------------------------
\section{Simulation Results}
\label{sec:results}

% Simulation setup and results will be added here once plots are generated.
% \begin{figure}[!t]
% \centering
% \includegraphics[width=3.5in]{vdc_response.png}
% \caption{DC bus voltage regulation under dynamic load and irradiance.}
% \label{fig:vdc}
% \end{figure}

Simulation results demonstrate:
\begin{itemize}
\item Stable DC bus voltage regulation to $300$\,V under step and ramp
      disturbances in both load and irradiance.
\item Near-MPPT PV power extraction across all operating conditions.
\item Smooth and coordinated activation of the electrolyzer during surplus
      and the fuel cell during deficit, governed by the EMS logic.
\item Effective battery SOC management within the $[0.20,\, 0.90]$ safe
      operating window.
\end{itemize}

%-------------------------------------------------
\section{Conclusion}
\label{sec:conclusion}

A distributed economic MPC framework for a PV--battery--hydrogen DC
microgrid has been presented. Detailed nonlinear models of all power
conversion stages, including the full single-diode PV array, alkaline
electrolyzer electrochemistry, and PEMFC polarization losses, were
incorporated directly into the local prediction models. The EMS layer
coordinates hydrogen system activation based on power surplus/deficit and
battery SOC. Simulation results confirm stable voltage regulation, MPPT
tracking, and coordinated energy dispatch under dynamic operating
conditions.

Future work includes experimental validation, real-time embedded
implementation, and comparison with centralized MPC benchmarks.

%-------------------------------------------------
\appendix
\section{System Parameters}
\label{app:params}

\begin{table}[!h]
\caption{DC Microgrid Component Parameters}
\label{tab:params}
\centering
\begin{tabular}{@{}lll@{}}
\toprule
Parameter & Symbol & Value \\ 
\midrule
\multicolumn{3}{l}{\textit{PV Array}} \\
Series modules          & $N_s$      & 5 \\
Parallel strings        & $N_{sh}$   & 20 \\
Cells per module        & $N_{cell}$ & 72 \\
Open-circuit voltage    & $V_{ocn}$  & 47.6 V \\
Short-circuit current   & $I_{scn}$  & 9.99 A \\
Series resistance       & $R_s$      & 0.20 $\Omega$ \\
Shunt resistance        & $R_{sh}$   & 250.0 $\Omega$ \\
Ideality factor         & $a$        & 1.30 \\
\midrule
\multicolumn{3}{l}{\textit{PV Boost Converter}} \\
Capacitance             & $C_s$      & 2000 $\mu$F \\
Inductance              & $L_s$      & 1 mH \\
\midrule
\multicolumn{3}{l}{\textit{Alkaline Electrolyzer}} \\
Series cells            & $N_{ae}$   & 30 \\
Cell area               & $A$        & 0.25 m$^2$ \\
Reversible voltage      & $U_{rev}$  & 1.229 V \\
\midrule
\multicolumn{3}{l}{\textit{PEMFC Stack}} \\
Series cells            & $N_{pe}$   & 300 \\
Operating temperature   & $T_{pe}$   & 323.15 K \\
Membrane area           & $A_{pe}$   & 50 cm$^2$ \\
\midrule
\multicolumn{3}{l}{\textit{Battery}} \\
Nominal voltage         & $V_{nom}$  & 240 V \\
Internal resistance     & $R_{int}$  & 0.15 $\Omega$ \\
Capacity                & $C_{Ah}$   & 100 Ah \\
\midrule
\multicolumn{3}{l}{\textit{Battery Bidirectional Converter}} \\
Inductance              & $L_b$      & 2 mH \\
\midrule
\multicolumn{3}{l}{\textit{DC Bus}} \\
Capacitance             & $C_{dc}$   & 30 mF \\
Voltage setpoint        & $V_{dc}^*$ & 300 V \\
\midrule
\multicolumn{3}{l}{\textit{DEMPC}} \\
Sampling interval       & $T_s$      & 40 $\mu$s \\
Prediction horizon      & $N_p$      & 2--3 \\
Reference horizon       & $N_s$      & 10 \\
\bottomrule
\end{tabular}
\end{table}

%-------------------------------------------------
\bibliographystyle{IEEEtran}
\bibliography{references}

\end{document}